% =====================================================================
%  5. Méthodes d'apprentissage
% =====================================================================
\clearpage
\section{Méthodes d'apprentissage}
\label{sec:methodes}

Nous comparons cinq familles supervisées vues en cours, choisies pour couvrir des
hypothèses différentes sur la forme de la relation entre conditions de soudage et
propriétés, puis une méthode semi-supervisée. Toutes sont implémentées avec
scikit-learn~\cite{pedregosa_2011} dans les notebooks 03 à 07. Dans la suite,
$\{(\bm{x}_i, y_i)\}_{i=1}^{n}$ désigne l'ensemble d'entraînement d'un pli, après
prétraitement.

\subsection{Méthodes supervisées}

\subsubsection{Référence : médiane de la cible}
Le modèle constant $\hat{y} = \operatorname{médiane}(y_1,\dots,y_n)$ ignore
complètement $\bm{x}$. Il fixe le \textbf{gain minimal} à dépasser : un modèle qui
ne fait pas mieux n'apprend rien d'utile des conditions de soudage. La médiane
minimise l'erreur absolue moyenne parmi les prédicteurs constants, ce qui en fait
la référence naturelle pour la MAE.

\subsubsection{Régression Ridge}
La régression Ridge~\cite{hoerl_1970} est une régression linéaire pénalisée :
\begin{equation}
  \hat{\bm{\beta}} = \arg\min_{\bm{\beta}}\; \sum_{i=1}^{n}\bigl(y_i - \beta_0 - \bm{x}_i^{\top}\bm{\beta}\bigr)^2 + \alpha\,\lVert\bm{\beta}\rVert_2^2 .
\end{equation}
La pénalité $\alpha$ réduit l'amplitude des coefficients et stabilise
l'estimation lorsque des entrées sont corrélées (courant, tension et apport
thermique, par exemple). Ridge sert de \textbf{contrôle linéaire} : si un modèle
non linéaire fait beaucoup mieux, la relation contient des non-linéarités ou des
interactions.

\subsubsection{Régression à vecteurs de support (SVR) à noyau RBF}
La SVR~\cite{smola_2004} cherche une fonction $f(\bm{x}) = \langle \bm{w}, \phi(\bm{x})\rangle + b$
aussi \og plate \fg{} que possible, en ignorant les erreurs inférieures à
$\varepsilon$ (perte $\varepsilon$-insensible) :
\begin{gather}
  \min_{\bm{w}, b}\; \tfrac{1}{2}\lVert\bm{w}\rVert^2 + C\sum_{i=1}^{n}\max\bigl(0,\, |y_i - f(\bm{x}_i)| - \varepsilon\bigr), \\
  k(\bm{x},\bm{x}') = \langle\phi(\bm{x}),\phi(\bm{x}')\rangle = \exp\bigl(-\gamma\lVert\bm{x}-\bm{x}'\rVert^2\bigr).
\end{gather}
Le noyau gaussien (RBF) permet des relations non linéaires lisses. $C$ règle le
compromis entre régularité et ajustement, $\gamma$ la portée du noyau et
$\varepsilon$ la largeur du tube de tolérance. Comme $\varepsilon$ s'exprime dans
l'unité de $y$, la cible est standardisée dans le train
(\texttt{TransformedTargetRegressor}) puis les prédictions sont reconverties.

\subsubsection{Forêt aléatoire}
Une forêt aléatoire~\cite{breiman_2001} moyenne $B$ arbres de régression, chacun
appris sur un échantillon bootstrap et choisissant ses coupures parmi un
sous-ensemble aléatoire de variables :
\begin{equation}
  \hat{y}(\bm{x}) = \frac{1}{B}\sum_{b=1}^{B} T_b(\bm{x}).
\end{equation}
Les arbres capturent les interactions et les effets de seuil (par exemple un
traitement thermique au-delà d'une certaine température) ; la moyenne réduit la
variance. La méthode est insensible aux transformations monotones des entrées et
robuste à des entrées peu informatives. Nous fixons $B = 300$ arbres et la
graine 42, et réglons la taille minimale des feuilles, la profondeur maximale et
la fraction de variables candidates à chaque coupure.

\subsubsection{$k$ plus proches voisins}
Le $k$NN~\cite{cover_1967} prédit la moyenne des cibles des $k$ voisins les plus
proches dans l'espace standardisé :
\begin{gather}
  \hat{y}(\bm{x}) = \frac{1}{k}\sum_{i \in \mathcal{N}_k(\bm{x})} y_i, \\
  d_p(\bm{x},\bm{x}') = \Bigl(\sum_j |x_j - x'_j|^p\Bigr)^{1/p}, \quad p = 1 \text{ (L1, Manhattan) ou } p = 2 \text{ (L2, euclidienne)}.
\end{gather}
C'est une méthode locale et sans hypothèse de forme. Elle est aussi le
\textbf{pendant supervisé de COREG}, qui repose sur deux $k$NN : la comparer
directement permet d'isoler l'effet des pseudo-labels.

\subsubsection{Grilles d'hyperparamètres}

Le \cref{tab:grilles} donne les grilles explorées. Elles ont été fixées
\emph{avant} les entraînements comme budgets raisonnables, et non optimisées a
posteriori. Au total, \textbf{39 configurations} sont évaluées par cible et par
pli externe.

\begin{table}[htbp]
\centering
\small
\caption{Familles supervisées, rôle dans la comparaison et grilles d'hyperparamètres.}
\label{tab:grilles}
\begin{tabularx}{\textwidth}{@{}l L l r@{}}
\toprule
Famille & Rôle & Grille & Config. \\
\midrule
Médiane de $y$ & Gain minimal à dépasser & aucune & 1 \\
Ridge & Contrôle linéaire régularisé & $\alpha \in \{0{,}01;\,0{,}1;\,1;\,10;\,100;\,1\,000\}$ & 6 \\
SVR RBF & Non-linéarité lisse à noyau & \makecell[l]{$C\in\{1,10,100\}$, $\gamma\in\{\text{scale}, 0{,}05\}$,\\ $\varepsilon\in\{0{,}05;\,0{,}2\}$} & 12 \\
Forêt aléatoire & Interactions et seuils & \makecell[l]{300 arbres ; feuille $\in\{1,2,5\}$,\\ profondeur $\in\{\text{libre}, 10\}$, \texttt{max\_features} $\in\{0{,}5;\,1\}$} & 12 \\
$k$NN & Méthode locale, contrôle de COREG & $k\in\{3,5,9,15\}$, $p\in\{1,2\}$, poids uniformes & 8 \\
\midrule
\multicolumn{3}{@{}l}{Total par cible et par pli externe} & 39 \\
\bottomrule
\end{tabularx}
\end{table}

\begin{decision}
Le critère de réglage est la \textbf{MAE interne moyenne} (\cref{subsec:metriques}).
Elle s'exprime dans l'unité de la cible et est moins dominée que la RMSE par
quelques séries très mal prédites. La famille elle-même peut être choisie dans la
boucle interne : on parle alors de \og sélection interne \fg{}, procédure dont le
score externe est une estimation honnête (\cref{sec:protocole}).
\end{decision}

\subsection{Apprentissage semi-supervisé : revue bibliographique}
\label{subsec:ssl_biblio}

L'apprentissage semi-supervisé cherche à améliorer un prédicteur en utilisant, en
plus des exemples étiquetés $\mathcal{D}_L$, des exemples non étiquetés
$\mathcal{D}_U$~\cite{chapelle_2006,vanengelen_2020}. Il n'aide que si la
distribution des entrées renseigne sur la fonction à apprendre, ce qui suppose
l'une des hypothèses suivantes :
\begin{itemize}
  \item \textbf{régularité} : deux entrées proches ont des sorties proches ;
  \item \textbf{amas} : les points d'un même amas partagent une sortie similaire ;
  \item \textbf{variété} : les données vivent près d'une variété de faible dimension.
\end{itemize}

Les principales familles de méthodes sont résumées au \cref{tab:ssl_familles}.

\begin{table}[htbp]
\centering
\small
\caption{Grandes familles de méthodes semi-supervisées et adéquation au projet.}
\label{tab:ssl_familles}
\begin{tabularx}{\textwidth}{@{}>{\raggedright\arraybackslash}p{3.2cm} L L@{}}
\toprule
Famille & Principe & Adéquation à notre problème \\
\midrule
Auto-apprentissage (\emph{self-training})~\cite{yarowsky_1995} & Un modèle étiquette ses prédictions les plus sûres et se réentraîne dessus & Simple, mais risque de renforcer ses propres erreurs ; la \og confiance \fg{} est mal définie en régression \\
Co-apprentissage (\emph{co-training})~\cite{blum_mitchell_1998} & Deux modèles sur deux vues indépendantes s'étiquettent mutuellement & Conçu pour la classification et pour deux vues naturelles, absentes ici \\
Méthodes à graphe (propagation de labels)~\cite{zhu_ghahramani_2002} & Les labels diffusent le long d'un graphe de similarité & Essentiellement transductives ; graphe sensible aux absences et aux échelles \\
Modèles génératifs, SVM transductives~\cite{chapelle_2006} & Modéliser $p(\bm{x}, y)$ ou placer la frontière dans les zones peu denses & Pensés surtout pour la classification \\
Régularisation par cohérence (réseaux profonds)~\cite{vanengelen_2020} & Prédictions stables sous perturbation & Exigent de grands volumes de données \\
\textbf{COREG}~\cite{zhou_li_2005} & Co-apprentissage pour la \textbf{régression}, deux $k$NN, confiance par réduction locale d'erreur & \textbf{Retenu} : conçu pour la régression, inductif, peu coûteux \\
\bottomrule
\end{tabularx}
\end{table}

\begin{decision}
Nous retenons \textbf{COREG} (\emph{CO-training REGressors}) de Zhou et
Li~\cite{zhou_li_2005}, pour cinq raisons :
\begin{enumerate}[leftmargin=1.5em]
  \item il est \textbf{conçu pour la régression}, alors que la plupart des méthodes le sont pour la classification ;
  \item il définit une \textbf{confiance mesurable} pour une prédiction numérique : la baisse d'erreur locale qu'elle provoque ;
  \item il n'exige pas deux vues naturelles indépendantes : la diversité provient de \textbf{deux métriques de distance} ;
  \item il repose sur le $k$NN, méthode vue en cours, ce qui permet un \textbf{contrôle supervisé exact} (mêmes $k$NN, mêmes labels, sans pseudo-labels) ;
  \item il est \textbf{inductif} : le modèle final prédit de nouvelles soudures sans réapprentissage, ce qui correspond à l'usage visé.
\end{enumerate}
\end{decision}

\subsection{L'algorithme COREG}
\label{subsec:coreg_algo}

COREG entraîne deux régresseurs $k$NN, $h_1$ (distance L1) et $h_2$ (distance L2),
sur les mêmes labels. À chaque tour, chaque régresseur examine des candidats du
pool $U$ et estime leur \textbf{confiance} : pour un candidat $\bm{x}_u$, on
calcule sa prédiction $\hat{y}_u = h_j(\bm{x}_u)$, on ajoute temporairement
$(\bm{x}_u, \hat{y}_u)$ au jeu de $h_j$ pour obtenir $h'_j$, puis on mesure la
réduction d'erreur quadratique sur les $k$ voisins étiquetés $\Omega_u$ de $\bm{x}_u$ :
\begin{equation}
  \Delta_u = \sum_{\bm{x}_i \in \Omega_u}\Bigl[\bigl(y_i - h_j(\bm{x}_i)\bigr)^2 - \bigl(y_i - h'_j(\bm{x}_i)\bigr)^2\Bigr].
  \label{eq:coreg_gain}
\end{equation}
Le candidat de plus grand gain positif est \textbf{transmis à l'autre régresseur},
qui l'ajoute à ses données comme pseudo-label. La prédiction finale est la
moyenne $\hat{y} = \tfrac{1}{2}\bigl(h_1(\bm{x}) + h_2(\bm{x})\bigr)$. La
\cref{fig:coreg} et l'\cref{algo:coreg} résument la procédure.

\begin{figure}[htbp]
\centering
\begin{tikzpicture}[
  reg/.style={rectangle, rounded corners=3pt, draw=csbordeaux, line width=0.7pt,
              fill=csgrisclair, align=center, text width=3.1cm, minimum height=1.3cm, font=\footnotesize},
  pool/.style={ellipse, draw=csgris, line width=0.6pt, fill=white, align=center,
               minimum width=3.2cm, minimum height=1.25cm, font=\footnotesize},
  fl/.style={-{Stealth[length=2mm]}, line width=0.8pt},
  lab/.style={font=\scriptsize, align=center}]
  \node[pool] (U) {Pool $U$ sans label\\ (train du pli seulement)};
  \node[reg, below left=0.9cm and 0.6cm of U] (h1) {$h_1$ : $k$NN, distance L1\\ labels $L$ + pseudo-labels de $h_2$};
  \node[reg, below right=0.9cm and 0.6cm of U] (h2) {$h_2$ : $k$NN, distance L2\\ labels $L$ + pseudo-labels de $h_1$};
  \node[reg, below=2.3cm of U, fill=white, draw=csgris, text width=3.6cm] (out) {Prédiction finale\\ $\hat{y} = \frac{1}{2}\bigl(h_1(\bm{x}) + h_2(\bm{x})\bigr)$};
  \draw[fl, color=csgris] (U) -- node[lab, left=2pt] {candidats\\ ($\leq 100$)} (h1);
  \draw[fl, color=csgris] (U) -- node[lab, right=2pt] {candidats\\ ($\leq 100$)} (h2);
  \draw[fl, color=orangeviz, bend left=12] (h1.east) to node[lab, above] {meilleur $\Delta_u>0$} (h2.west);
  \draw[fl, color=bleuviz, bend left=12] (h2.west) to node[lab, below] {meilleur $\Delta_u>0$} (h1.east);
  \draw[fl, color=csgris] (h1.south) |- (out.west);
  \draw[fl, color=csgris] (h2.south) |- (out.east);
\end{tikzpicture}
\caption{Principe de COREG. Chaque $k$NN choisit dans le pool l'exemple dont sa
propre prédiction réduit le plus l'erreur locale, puis le transmet à l'autre
$k$NN comme pseudo-label.}
\label{fig:coreg}
\end{figure}

\begin{algorithm}[htbp]
\caption{COREG tel qu'implémenté dans le notebook 07 (un pli, une cible)}
\label{algo:coreg}
\small
\begin{algorithmic}[1]
\Require labels du train $L$ (prétraités), pool $U$ du train, $k$, $T_{\max}=100$, $N_{\text{cand}}=100$
\State Ajuster le prétraitement sur $L$ ; le figer ; transformer $U$
\State Dédupliquer $U$ par représentation exacte ; retirer les $\bm{x}$ déjà présents dans $L$
\State $L_1 \gets L$, $L_2 \gets L$ ; $h_1 \gets k\text{NN}_{L1}(L_1)$, $h_2 \gets k\text{NN}_{L2}(L_2)$
\For{$t = 1, \dots, T_{\max}$}
  \State Tirer (graine dérivée du pli) au plus $N_{\text{cand}}$ candidats actifs $U' \subseteq U$
  \For{$j \in \{1, 2\}$}
    \State Pour chaque $\bm{x}_u \in U'$ : $\hat{y}_u \gets h_j(\bm{x}_u)$ et calculer $\Delta_u$ (\cref{eq:coreg_gain})
    \State $u_j^\star \gets \arg\max_u \Delta_u$ ; retenir si $\Delta_{u_j^\star} > \tau = 10^{-12}\max(1, \operatorname{var}(y))\,k$
    \State Retirer $\bm{x}_{u_j^\star}$ de $U'$ et de $U$
  \EndFor
  \State $L_2 \gets L_2 \cup \{(\bm{x}_{u_1^\star}, \hat{y}_{u_1^\star})\}$, $L_1 \gets L_1 \cup \{(\bm{x}_{u_2^\star}, \hat{y}_{u_2^\star})\}$ ; mettre à jour $h_1, h_2$
  \If{aucun ajout accepté \textbf{ou} $U$ épuisé}
    \State \textbf{arrêter}
  \EndIf
\EndFor
\State \Return $\hat{y}(\bm{x}) = \frac{1}{2}\bigl(h_1(\bm{x}) + h_2(\bm{x})\bigr)$
\end{algorithmic}
\end{algorithm}

Le \cref{tab:coreg_choix} distingue ce qui vient de l'article original et ce
qui relève de nos choix d'implémentation, tous fixés \emph{avant} le calcul.

\begin{table}[htbp]
\centering
\small
\caption{Paramétrage de COREG et écarts à l'article original.}
\label{tab:coreg_choix}
\begin{tabularx}{\textwidth}{@{}>{\raggedright\arraybackslash}p{3.4cm} L@{}}
\toprule
Élément & Choix du projet \\
\midrule
Régresseurs & Deux $k$NN uniformes sur les mêmes 16 entrées $+$ procédé ; distances L1 et L2 (l'article utilise des distances de Minkowski d'ordres 2 et 5) \\
Nombre de voisins & Même $k$ pour les deux, choisi dans $\{3, 5, 9, 15\}$ par MAE interne \\
Prétraitement & Appris sur les vrais labels du train, figé pendant toute la boucle \\
Confiance & \cref{eq:coreg_gain} ; voisinage dans le jeu courant, pseudo-labels précédents inclus ; mise à jour exacte vérifiée contre un réapprentissage explicite \\
Candidats et tours & Au plus 100 candidats actifs par tour, au plus 100 tours, un ajout par régresseur et par tour, aucun identifiant ajouté deux fois \\
Pool & Déduplication des représentations identiques (plus petit identifiant conservé) ; retrait des $\bm{x}$ déjà présents dans les vrais labels \\
Arrêt & Pool épuisé, aucun gain positif dans le pool actif, ou 100 tours ; jamais sur un label de validation \\
Égalités & Ordre stable : vrais labels triés par identifiant, puis représentants du pool \\
\bottomrule
\end{tabularx}
\end{table}

\begin{precaution}
Un pseudo-label est une \emph{estimation}, pas une mesure mécanique. Le gain
local $\Delta_u$ mesure la cohérence d'une prédiction avec ses voisins étiquetés ;
il ne garantit pas qu'elle soit exacte, et deux $k$NN qui partagent les mêmes
entrées peuvent propager une même erreur.
\end{precaution}

\begin{aretenir}
Cinq familles supervisées couvrent le constant, le linéaire, le non-linéaire lisse,
les arbres et le local ; 39 configurations sont réglées par MAE interne.
COREG est retenu comme méthode semi-supervisée parce qu'il est conçu pour la
régression, inductif et directement comparable à un contrôle $k$NN supervisé.
\end{aretenir}
